% Sequential recommendation aims to suggest relevant items based on the user's historical behaviors. Many algorithms have been deployed to better understand the sequential patterns of user's behaviors. Markov-chain based methods \cite{SIMILARITYMC, MCAULEYMC2016, MCAULEYMC2017, he2016fusing} use the traditional Markov Chain model by assuming that users' behaviors are affected only by their last few behaviors. FPMC \cite{rendle2010factorizing} merges Markov-chains with matrix factorization method for next-basket recommendation. Since the first suggestion by GRU4Rec\cite{GRU4Rec}, many RNN-based methods \cite{hidasi2018recurrent, RNN4Rec, HIERARCHICALGRU4REC, DREAM, wu2017recurrent, wu2017sequential, DINEVOLUTION} sought to bring the success of RNN in language understanding into item sequence understanding. To overcome the limitations of RNN models (namely, the strong order constraint imposed by RNN), CNN-based method CASER \cite{CASER}, and others\cite{CNNGENERATIVE, HIERARCHICALCNN} proposed to use CNNs to better capture the dependencies between items. Some works adopted GNN to understand user's session as a graph\cite{SESSIONGNN}.

Sequential recommendation aims to suggest relevant items based on the user's sequential history. Markov-chain based methods~\cite{SIMILARITYMC, MCAULEYMC2016, TRANSREC, he2016fusing} assume that users' behaviors are affected only by their last few behaviors. FPMC~\cite{rendle2010factorizing} merges Markov-chains with matrix factorization method for next-basket recommendation. Since the first suggestion by GRU4Rec~\cite{GRU4Rec}, many RNN-based methods~\cite{hidasi2018recurrent, RNN4Rec, HIERARCHICALGRU4REC, DREAM, wu2017recurrent, wu2017sequential, DINEVOLUTION} brought the success of RNN into item sequence understanding. To overcome the strong order constraint of RNN models, CNN-based methods~\cite{CASER, CNNGENERATIVE, HIERARCHICALCNN} were proposed. Some works adopted graph neural network (GNN) to understand user's session as a graph~\cite{SESSIONGNN}.

%As for attention mechanisms, NARM \cite{NARM}, STAMP \cite{STAMP}, etc incorporated vanilla attention mechanism with RNN. More recently, SASRec\cite{SASRec} successfully employed self-attention mechanism, which was a huge success in NLP areas\cite{TRANSFORMER, BERT, TRANSFORMER-XL, XLNET}. BERT4Rec \cite{BERT4Rec} improved SASRec by replacing their self-attention blocks with bi-directional transformers and using Cloze task to train the model. TiSASRec \cite{TISASRec} enhanced SASRec by merging timestamp information into self-attention operations. Our work seeks to further expand this trend by suggesting a novel model architecture that uses multiple types of positional embeddings and self-attention operations to better capture the diverse patterns in user's behaviors.

As for attention mechanisms, NARM~\cite{NARM} and STAMP~\cite{STAMP} incorporated vanilla attention mechanism with RNN. More recently, SASRec~\cite{SASRec} successfully employed self-attention mechanism, which was a huge success in NLP areas~\cite{TRANSFORMER, BERT, TRANSFORMER-XL, XLNET}. BERT4Rec~\cite{BERT4Rec} improved SASRec by adopting Transformer~\cite{TRANSFORMER} and Cloze-task based training method. TiSASRec~\cite{TISASRec} enhanced SASRec by merging timestamp information into self-attention operations. Our work seeks to further expand this trend by suggesting a novel model architecture that uses multiple types of temporal embeddings  and self-attention operations to better capture the diverse patterns in user's behaviors.